\documentclass{article}
\usepackage[T1]{fontenc}
\usepackage{lmodern}
\usepackage{graphicx}
\usepackage{amsmath,amssymb}
\usepackage[a4paper,margin=2cm]{geometry}
\usepackage[absolute,overlay]{textpos}
\setlength{\TPHorizModule}{1mm}
\setlength{\TPVertModule}{1mm}
\usepackage[indonesian]{babel}
\usepackage{hyperref}
\setlength{\emergencystretch}{2em}
% unit: Halaman 9

\begin{document}

% === Halaman 9 (python/native) ===

9

◼

Ini sering terjadi jika sebagian isi pesan bisa ditebak (misalnya header file, 

protokol standar, salam pembuka email).

◼

Beberapa pesan yang formatnya terstruktur membuka peluang untuk menerka 

plainteks dari cipherteks yang bersesuaian. 

      Contoh: 


From dan To  di dalam e-mail, 

           ”Dengan hormat”, wassalam, pada surat resmi.

            \#include, program, di dalam source code

◼Dengan menggunakan pasangan plainteks dan cipherteks, kriptanalis dapat 

menemukan kunci enkripsi  (lihat pembahasan cipher klasik affine cipher dan 

hill cipher)

\end{document}
